\section{Introduction}
TODO: Introduce the Bellman Operator and Bellman Optimality Operator.
The \textbf{Bellman equation} is the heart of reinforcement learning.
All reinforcement learning algorithms are designed to solve this equation.
\section{Bellman Equation}
    \subsection{State-value Function}
        Before diving into the Bellman equation, we introduce an important concept called the \textbf{value of a state}.
        \begin{defn}{State Value}{state-value}
            Let the agent follows the policy $\pi$. The \textbf{state value} $v_\pi(s)$ of state $s$ is defined as the expected return:
            \begin{equation}
                v_\pi(s) = \E[G_t|S_t=s] \label{eq:state-value}
            \end{equation}
            where $G_t$ is the discounted return $G_t=\sum_{k=t+1}^{\infty}\gamma^{k-t-1}R_{k}$.
            The function $v_\pi$ itself is called \textbf{state-value function}.
        \end{defn}
        State values are used to evaluate the policy $\pi$. Policies that generate greater state values are better. Then, how can we
        obtain te state values? The answer is to solve the Bellman equation.
    \subsection{Bellman Equation for State-Value Function}
        The Bellman equation is simply a set of linear equations that describe the relationships of each state values.
        Bellman equation can be obtained by using the recursive property of return:
        \begin{align}
            G_t &= R_{t+1} + \gamma R_{t+2} + \gamma^2 R_{t+3} + \gamma^3 R_{t+4} + \cdots \\
                &= R_{t+1} + \gamma (R_{t+2} + \gamma R_{t+3} + \gamma^2 R_{t+4}) + \cdots \\
                &= R_{t+1} + \gamma G_{t+1}.
        \end{align}
        Substituting equation \ref{eq:state-value} with the recursive realationship, we obtain
        \begin{align}
            v_\pi(s) &= \E[G_t|S_t=s] \\
                     &= \E[R_{t+1} + \gamma G_{t+1}|S_t=s] \\
                     &= \E[R_{t+1}|S_t=s] + \gamma \E[G_{t+1}|S_t=s]
        \end{align} 
        Using the law of total expectation \ref{eq:cond-LTE}, we get
        \begin{align}
            \E[R_{t+1}|S_t=s] &= \sum_{a \in \mathcal{A}(s)}\pi(a|s)\E[R_{t+1}|S_t=s,A_t=a] \\
                              &= \sum_{a \in \mathcal{A}(s)}\pi(a|s)\sum_{r \in \mathcal{R}(s,a)}p(r|s,a)r
        \end{align}
        and
        \begin{align}
            \E[G_{t+1}|S_t=s] &= \sum_{s' \in \mathcal{S}}\E[G_{t+1}|S_{t+1}=s',S_t=s]p(s'|s) \\
                              &= \sum_{s' \in \mathcal{S}}\E[G_{t+1}|S_{t+1}=s']p(s'|s) \quad \text{(Markov property)} \\
                              &= \sum_{s' \in \mathcal{S}}v_\pi(s')p(s'|s) \\
                              &= \sum_{s' \in \mathcal{S}}v_\pi(s')\sum_{a \in \mathcal{A}(s)}p(s'|s,a).
        \end{align}
        The final result is given in the following box:
        \begin{thm}{Bellman Equation for State-Value Function}{BellmanEq-SV}
            Let $\pi$ be a policy, and let $v_\pi$ be a state-value function. Then, the following Bellman equation is satisfied:
            \begin{equation}
                v_\pi(s) = \sum_{a \in \mathcal{A}(s)}\pi(a|s)\sum_{s' \in \mathcal{S}}\sum_{r \in \mathcal{R}(s,a)}p(s',r|s,a)\big[ r + \gamma v_\pi(s') \big]
            \end{equation}
        \end{thm}
    \subsection{Matrix-Vector form of Bellman Equation}
        We can reformulate Bellman equation in terms of matrix-vector form. For this, we rewrite the Bellman equation given in Theorem \ref{thm:BellmanEq-SV} as
        \begin{equation}
            v_\pi(s) = r_\pi(s) + \gamma\sum_{s' \in \mathcal{S}}p_\pi(s'|s)v_\pi(s'),
        \end{equation}
        where
        \begin{align}
            r_\pi(s) &= \sum_{a \in \mathcal{A}(s)}\pi(a|s)\sum_{r \in \mathcal{R}(s,a)}p(r|s,a)r \\
            p_\pi(s'|s) &= \sum_{a \in \mathcal{A}(s)}\pi(a|s)p(s'|s,a).
        \end{align}
        The $r_\pi$ and $p_\pi$ denotes the mean of immediate rewards and probability of transitioning from $s$ to $s'$ under policy $\pi$, respectively.
        Suppose that we indexed all states as $s_i$ with $i = 1, \dots, n$, where $n=|\mathcal{S}|$.
        Let
        \begin{equation}
            v_\pi = \begin{bmatrix}
                    v_\pi(s_1) \\
                    \vdots \\
                    v_\pi(s_n)
                    \end{bmatrix},
            r_\pi = \begin{bmatrix}
                    r_\pi(s_1) \\
                    \vdots \\
                    r_\pi(s_n)
                    \end{bmatrix},
            P_\pi = \begin{bmatrix}
                    p_\pi(s_1|s_1) & p_\pi(s_2|s_1) & \dots & p_\pi(s_n|s_1) \\
                    p_\pi(s_1|s_2) & p_\pi(s_2|s_2) & \dots & p_\pi(s_n|s_2) \\
                    \vdots & \vdots &  & \vdots \\
                    p_\pi(s_1|s_n) & p_\pi(s_2|s_n) & \dots & p_\pi(s_n|s_n)
                    \end{bmatrix}. \label{eq:VecForm}
        \end{equation}
        Then, Theorem \ref{thm:BellmanEq-SV} can be written in the following form:
        \begin{thm}{Matrix-Vector form of Bellman Equation}{MVFormBellmanEq}
        Let $v_\pi$, $r_\pi$ and $P_\pi$ be given as \ref{eq:VecForm}. Then, following is satisfied:
            \begin{equation}
                v_\pi = r_\pi + \gamma P_\pi v_\pi
            \end{equation}
        \end{thm}
        \begin{rem}{Properties of Matrix $P_\pi$}{ProbMatProp}
            Note that the matrix $P_\pi$ given in equation \ref{eq:VecForm} satisfies following properties:
            \begin{itemize}
                \item All the elements are nonnegative;
                \item The sum of the values of every row equals 1. Matrix with this property is called a \emph{stochastic matrix}.
            \end{itemize}
        \end{rem}
\section{Solving The Bellman Equation}
    There are two methods in solving the Bellman equation. The closed-form solution, and the iterative solution.
    \subsection{Closed-form Solution}
        Since we have represented the Bellman equation as the matrix form, solving it is very easy.
        \begin{thm}{Closed-form Solution of The Bellman Equation}{ClosedFormSolBellmanEq}
            Let $\pi$ be a given policy, and $v_\pi$, $r_\pi$, $P_\pi$ be given as \ref{eq:VecForm}.
            The solution to the equation $v_\pi = r_\pi + \gamma P_\pi v_\pi$ is given as
            \begin{equation}
                v_\pi = (I-\gamma P_\pi)^{-1}r_\pi
            \end{equation}
        \end{thm}
        \begin{proof}
            Manipulate the equation as following:
            \begin{align}
                &v_\pi = r_\pi + \gamma P_\pi v_\pi \\
                &v_\pi - \gamma P_\pi v_\pi = r_\pi \\
                &(I - \gamma P_\pi)v_\pi = r_\pi.
            \end{align}
            Now, what we have to show is that the matrix $I - \gamma P_\pi$ is invertible.
            We use the fact that $A$ is invertible $\Leftrightarrow$ eigenvalues of $A$ are nonzero.
            To show this, we use Geshgorin Circle Theorem \ref{thm:Geshgorin}. Note that the $i$-th Geshgorin circles has a center at $1-\gamma p(s_i|s_i)$, with
            radius $\sum_{j \neq i}\gamma p_\pi(s_j|s_i)$. Since $P_\pi$ is stochastic matrix, we know that $\sum_{j}\gamma p(s_j|s_i)=\gamma < 1$.
            That is, $\sum_{j \neq i}\gamma p_\pi(s_j|s_i) < 1 - \gamma p_\pi(s_i|s_i)$, which means that the $i$-th Geshgorin circle cannot contain the origin.
            According to the Geshgorin Circle Theorem, all the eigenvalues cannot be zero.
        \end{proof}
        TODO: add properties of $I - \gamma P_\pi$.
    \subsection{Iterative Solution}
        The direct method for solving the Bellman equation is not used very well in practice, since the inverse matrix computation is very costly.
        Instead of the close-form method, in practice, we use iterative method to solve the Bellman equation.
        \begin{thm}{Iterative Solution of The Bellman Equation}{IterativeSolBellmanEq}
            Let $\pi$ be a policy and $v_\pi$, $r_\pi$, and $P_\pi$ be that of \ref{eq:VecForm}. Then, the
            solution of \ref{thm:MVFormBellmanEq} equals the limit of following sequence $\{v_k\}$:
            \begin{equation}
                v_{k+1}=r_\pi + \gamma P_\pi v_k \label{eq:IterSolBellmanEq}
            \end{equation}
        \end{thm}
        \begin{proof}
            What we have to show is that $v_k \to v_\pi$ as $k \to \infty$. For this, we let $\delta_k = v_k - v_\pi$ and show that $\delta_k \to 0$.
            Substituting $v_{k+1} = \delta_{k+1} + v_\pi$ and $v_k = \delta_k + v_\pi$ into equation \ref{eq:IterSolBellmanEq},
            we have
            \begin{equation}
                \delta_{k+1} = \gamma P_\pi \delta_k.
            \end{equation}
            Since $P_\pi$ is a stochastic matrix and $0 < \gamma < 1$, we know that $\delta_k \to 0$ as $k \to \infty$.
        \end{proof}

\section{Bellman Equation for Action-Value Function}
    Now we introduce the \textbf{action value}, which denotes the value of an action at a state.
    \subsection{Action Value}
        The action value is defined as following:
        \begin{defn}{Action Value}{action-value}
            Let $\pi$ be a given policy. Then, the \textbf{action value} of a state-action pair is defined as
            \begin{equation}
                q_\pi(s,a) = \E[G_t|S_t=s,A_t=a].
            \end{equation}
        \end{defn}
        \begin{rem}{Properties of Action Value}{ActionValProp}
            Action value $q_\pi(s,a)$ has following properties.
            \begin{itemize}
                \item $v_\pi=\sum_{a \in \mathcal{A}(s)} \pi(a|s)q_\pi(s,a)$
                \item $q_\pi(s,a) = \sum_{r \in \mathcal{R}(s,a)}p(r|s,a)r + \gamma \sum_{s' \in \mathcal{S}}p(s'|s,a)v_\pi(s')$
            \end{itemize}
        \end{rem}
        \begin{proof}
            \begin{itemize}
                \item Use the law of total expectation \ref{eq:cond-LTE}.
                \item Use the result above and \ref{thm:BellmanEq-SV}.
            \end{itemize}
        \end{proof}
    \subsection{Bellman Equation for Action-Value Function}
        The Bellman equation for action-value function is already in our hand, since actually the Remark \ref{rem:ActionValProp}
        is what all we need.
        \begin{thm}{Bellman Equation for Action-value Function}{BellmanEq-AV}
            Let $\pi$ be a given policy, and let $q_\pi(s,a)$ be an action value. Then, the following Bellman equation is satisfied:
            \begin{align}
                q_\pi(s,a) &= \sum_{r \in \mathcal{R}(s,a)}p(r|s,a)r + \gamma \sum_{s' \in \mathcal{S}}p(s'|s,a)\sum_{a' \in \mathcal{A}(s')}\pi(a'|s')q_\pi(s',a') \\
                           &= \sum_{r \in \mathcal{R}(s,a)}p(r|s,a)r + \gamma \sum_{s' \in \mathcal{S}}p(s'|s,a)v_\pi(s') 
            \end{align}
            Moreover, the matrix-vector form is given as
            \begin{equation}
                q_\pi = \tilde{r} + \gamma \tilde{P} \Pi q_\pi
            \end{equation}
            where $[q_\pi]_{(s,a)}=q_\pi(s,a)$, $\tilde{r}_{(s,a)}=\sum_{r \in \mathcal{R}(s,a)}p(r|s,a)r$, $[\tilde{P}]_{(s,a),s'}=p(s'|s,a)$ and $\Pi$ is a block diagonal matrix
            such that $\Pi_{s',(s',a')}=\pi(s'|a')$ for $a' \in \mathcal{A}(s')$ and other entries are all zero.
        \end{thm}
        \begin{proof}
            Use the Remark \ref{rem:ActionValProp}. For the matrix-vector form, try to derive it, as we did in Theorem \ref{thm:MVFormBellmanEq}.
        \end{proof}